\textbf{Phần 1. Trắc nghiệm lựa chọn}
\setcounterpageref{ex}{0}

\begin{ex} %[4D13N70]
Kết quả của phép tính $\dfrac{1}{4} + \dfrac{2}{4}$ là
\choice
{$\dfrac{3}{8}$}
{\True $\dfrac{3}{4}$}
{$\dfrac{2}{4}$}
{$\dfrac{1}{2}$}
\loigiai{
$\dfrac{1 + 2}{4} = \dfrac{3}{4}$.\\
\textbf{Đáp án của câu hỏi là B}
}
\end{ex}

\begin{ex} %[4D13N70]
Kết quả của phép tính $\dfrac{3}{5} \times \dfrac{2}{7}$ là
\choice
{\True $\dfrac{6}{35}$}
{$\dfrac{5}{12}$}
{$\dfrac{6}{12}$}
{$\dfrac{5}{35}$}
\loigiai{
$\dfrac{3 \times 2}{5 \times 7} = \dfrac{6}{35}$.\\
\textbf{Đáp án của câu hỏi là A}
}
\end{ex}

\begin{ex} %[4D13H70]
Kết quả của phép trừ $\dfrac{5}{6} - \dfrac{1}{2}$ là
\choice
{$\dfrac{4}{4}$}
{\True $\dfrac{1}{3}$}
{$\dfrac{2}{3}$}
{$\dfrac{1}{6}$}
\loigiai{
$\dfrac{5}{6} - \dfrac{3}{6} = \dfrac{2}{6} = \dfrac{1}{3}$.\\
\textbf{Đáp án của câu hỏi là B}
}
\end{ex}

\begin{ex} %[4D13H70]
Kết quả của phép chia $\dfrac{4}{9} : \dfrac{2}{3}$ là
\choice
{$\dfrac{8}{27}$}
{$\dfrac{3}{2}$}
{$\dfrac{1}{3}$}
{\True $\dfrac{2}{3}$}
\loigiai{
$\dfrac{4}{9} \times \dfrac{3}{2} = \dfrac{2}{3}$.\\
\textbf{Đáp án của câu hỏi là D}
}
\end{ex}

\begin{ex} %[4D13H70]
$\dfrac{2}{3}$ của $30\text{ kg}$ là
\choice
{$15\text{ kg}$}
{\True $20\text{ kg}$}
{$25\text{ kg}$}
{$10\text{ kg}$}
\loigiai{
$30 \times \dfrac{2}{3} = 20\text{ kg}$.\\
\textbf{Đáp án của câu hỏi là B}
}
\end{ex}

\begin{ex} %[4D13H70]
Tính thuận tiện: $\dfrac{3}{5} \times \dfrac{7}{9} + \dfrac{3}{5} \times \dfrac{2}{9}$.
\choice
{\True $\dfrac{3}{5}$}
{$1$}
{$\dfrac{9}{9}$}
{$\dfrac{6}{5}$}
\loigiai{
$\dfrac{3}{5} \times \left(\dfrac{7}{9} + \dfrac{2}{9}\right) = \dfrac{3}{5} \times 1 = \dfrac{3}{5}$.\\
\textbf{Đáp án của câu hỏi là A}
}
\end{ex}

\begin{ex} %[4D13V70]
Một mảnh đất hình chữ nhật có chiều dài $\dfrac{4}{5}\text{ m}$, chiều rộng $\dfrac{1}{2}\text{ m}$. Diện tích mảnh đất là
\choice
{$\dfrac{5}{7}\text{ m}^2$}
{$\dfrac{4}{10}\text{ m}^2$}
{\True $\dfrac{2}{5}\text{ m}^2$}
{$\dfrac{3}{10}\text{ m}^2$}
\loigiai{
$\dfrac{4}{5} \times \dfrac{1}{2} = \dfrac{4}{10} = \dfrac{2}{5}\text{ m}^2$.\\
\textbf{Đáp án của câu hỏi là C}
}
\end{ex}

\begin{ex} %[4D13V70]
Một người có $60\text{ kg}$ gạo, đã bán $\dfrac{3}{4}$ số gạo đó. Số gạo còn lại là
\choice
{\True $15\text{ kg}$}
{$45\text{ kg}$}
{$20\text{ kg}$}
{$30\text{ kg}$}
\loigiai{
Đã bán: $60 \times \dfrac{3}{4} = 45\text{ kg}$. Còn lại: $60 - 45 = 15\text{ kg}$.\\
\textbf{Đáp án của câu hỏi là A}
}
\end{ex}

\begin{ex} %[4D13V70]
Tìm $x$ biết: $x - \dfrac{1}{3} = \dfrac{2}{5}$.
\choice
{$\dfrac{1}{2}$}
{$\dfrac{3}{8}$}
{$\dfrac{1}{15}$}
{\True $\dfrac{11}{15}$}
\loigiai{
$x = \dfrac{2}{5} + \dfrac{1}{3} = \dfrac{6 + 5}{15} = \dfrac{11}{15}$.\\
\textbf{Đáp án của câu hỏi là D}
}
\end{ex}

\begin{ex} %[4D13V70]
Tìm $x$ biết: $x : \dfrac{3}{4} = \dfrac{8}{9}$.
\choice
{$\dfrac{32}{27}$}
{$\dfrac{5}{12}$}
{\True $\dfrac{2}{3}$}
{$\dfrac{6}{12}$}
\loigiai{
$x = \dfrac{8}{9} \times \dfrac{3}{4} = \dfrac{24}{36} = \dfrac{2}{3}$.\\
\textbf{Đáp án của câu hỏi là C}
}
\end{ex}

\begin{ex} %[4D13V70]
Tính giá trị biểu thức: $\left(\dfrac{1}{2} + \dfrac{1}{3}\right) \times \dfrac{6}{5}$.
\choice
{$\dfrac{5}{6}$}
{$\dfrac{6}{5}$}
{$2$}
{\True $1$}
\loigiai{
$\dfrac{5}{6} \times \dfrac{6}{5} = 1$.\\
\textbf{Đáp án của câu hỏi là D}
}
\end{ex}

\begin{ex} %[4D13C70]
Tính giá trị biểu thức: $A = \dfrac{1}{2} \times \dfrac{2}{3} \times \dfrac{3}{4} \times \dots \times \dfrac{19}{20}$.
\choice
{$\dfrac{1}{10}$}
{$\dfrac{1}{19}$}
{\True $\dfrac{1}{20}$}
{$\dfrac{2}{20}$}
\loigiai{
$A = \dfrac{1}{20}$.\\
\textbf{Đáp án của câu hỏi là C}
}
\end{ex}

\textbf{Phần 2. Trắc nghiệm đúng sai}
\setcounterpageref{ex}{0}

\begin{ex} %[4D13H70]
Xét tính đúng sai của các phép toán phân số:
\choiceTF
{\True $\dfrac{1}{3} + \dfrac{2}{3} = 1$.}
{\True $\dfrac{3}{4} \times \dfrac{4}{3} = 1$.}
{$\dfrac{5}{6} - \dfrac{1}{6} = \dfrac{4}{0}$.}
{\True $\dfrac{5}{6} - \dfrac{1}{6} = \dfrac{4}{6} = \dfrac{2}{3}$.}
\loigiai{
\begin{itemize}
    \item $1/3 + 2/3 = 1$. Đúng.
    \item $3/4 \times 4/3 = 1$. Đúng.
    \item Mẫu số không thể bằng 0. Sai.
    \item $5/6 - 1/6 = 2/3$. Đúng.
\end{itemize}
\begin{itemize}
    \item a) \textbf{Đúng}.
    \item b) \textbf{Đúng}.
    \item c) \textbf{Sai}.
    \item d) \textbf{Đúng}.
\end{itemize}
}
\end{ex}

\begin{ex} %[4D13V70]
Một thùng có $24\text{ l}$ dầu: Lần 1 lấy $\dfrac{1}{3}$ thùng, lần 2 lấy $\dfrac{1}{4}$ thùng. Xét tính đúng sai:
\choiceTF
{\True Lần 1 lấy được $8\text{ l}$ dầu.}
{\True Lần 2 lấy được $6\text{ l}$ dầu.}
{Cả hai lần lấy được tất cả $15\text{ l}$ dầu.}
{\True Trong thùng còn lại $10\text{ l}$ dầu.}
\loigiai{
\begin{itemize}
    \item Lần 1: $24 \times 1/3 = 8\text{ l}$. Đúng.
    \item Lần 2: $24 \times 1/4 = 6\text{ l}$. Đúng.
    \item Tổng lấy: $8 + 6 = 14\text{ l} \neq 15\text{ l}$. Ý 3 Sai, Ý 4 Đúng ($24 - 14 = 10\text{ l}$).
\end{itemize}
\begin{itemize}
    \item a) \textbf{Đúng}.
    \item b) \textbf{Đúng}.
    \item c) \textbf{Sai}.
    \item d) \textbf{Đúng}.
\end{itemize}
}
\end{ex}

\textbf{Phần 3. Trả lời ngắn}
\setcounterpageref{ex}{0}

\begin{ex} %[4D13H70]
Tính tổng: $\dfrac{2}{5} + \dfrac{1}{2}$ (dạng a/b).
\shortans{9/10}
\loigiai{
\textbf{Đáp số:} $9/10$\\
$4/10 + 5/10 = 9/10$.
}
\end{ex}

\begin{ex} %[4D13H70]
Tính hiệu: $\dfrac{7}{8} - \dfrac{3}{4}$ (dạng a/b).
\shortans{1/8}
\loigiai{
\textbf{Đáp số:} $1/8$\\
$7/8 - 6/8 = 1/8$.
}
\end{ex}

\begin{ex} %[4D13V70]
Tính tích: $\dfrac{5}{9} \times \dfrac{3}{10}$ (rút gọn tối giản dạng a/b).
\shortans{1/6}
\loigiai{
\textbf{Đáp số:} $1/6$\\
$15/90 = 1/6$.
}
\end{ex}

\begin{ex} %[4D13V70]
Tính thương: $\dfrac{3}{5} : \dfrac{9}{10}$ (rút gọn tối giản dạng a/b).
\shortans{2/3}
\loigiai{
\textbf{Đáp số:} $2/3$\\
$3/5 \times 10/9 = 2/3$.
}
\end{ex}

\begin{ex} %[4D13V70]
Tìm $\dfrac{3}{7}$ của $42$.
\shortans{18}
\loigiai{
\textbf{Đáp số:} $18$\\
$42 \times 3/7 = 18$.
}
\end{ex}

\begin{ex} %[4D13C70]
Tính nhanh: $\dfrac{3}{4} \times \dfrac{5}{7} + \dfrac{3}{4} \times \dfrac{2}{7}$ (dạng a/b).
\shortans{3/4}
\loigiai{
\textbf{Đáp số:} $3/4$\\
$3/4 \times 1 = 3/4$.
}
\end{ex}
